%% P13 --- Graphs, DAGs and random walks
%%
%% Every number in this program that the reader cannot do in their head comes
%% from code/p13_graphs.py and is pulled in with \val{}. The console listing is
%% a committed transcript written by that script.
%%
%% THE THESIS: a graph is a set plus a relation; the adjacency matrix turns
%% every question about paths into a question about matrix powers; and a DAG's
%% topological order is what makes evaluation well defined -- so an agent
%% workflow, a build system and a neural network are one object.
%%
%% WHAT THIS PROGRAM IS OWED, read out of the files rather than remembered:
%%   F10  DEFERS "a graph as a set plus a relation" here BY NAME, in its own
%%        file header, and gives the sets, the union rule and the pair count.
%%   P12  owns the counting arguments and deliberately did not use them on
%%        graphs. Its documents-and-pairs example continues here.
%%   P06  gives the theorem section 3 rests on: a matrix product IS a
%%        composition, and row-meets-column falls out of that. Which is
%%        exactly why A^k counts walks of length k.
%%   P10  gives eigenvectors. A stationary distribution is an eigenvector for
%%        eigenvalue 1, which is the same calculation under another name.
%%   P03  gives the cost vocabulary for section 2.
%%
%% WHAT THIS PROGRAM OWES FORWARD: P16 declares P13 as a dependency and needs
%% the computation graph defined properly, one program before autodiff uses it.
%%
%% WHAT IT LEAVES ALONE, checked against tools/programs.json:
%%   the Jacobian, forward and reverse mode, what backward() computes -> P16
%%   quantifiers, induction, writing a proof                          -> P14
%%   what a probability IS, and conditioning                          -> P23
%%
%% Twin: programs/pl/P13-graphs-dags.tex. Frame for frame, macro for macro, number
%% for number; tools/parity.py compares them.

\program{Graphs, DAGs and random walks}
\label{prog:P13}
\index{graph}
\index{graph!adjacency matrix}

\begin{outcomes}
\outcome{1--8}{say what a graph is with no picture, and count its edges
  against the pairs they are drawn from}
\outcome{9--13}{choose between an adjacency matrix and an adjacency list from
  the density and the operation, with the bill}
\outcome{14--20}{read a power of the adjacency matrix as a count of walks, and
  say what that makes the depth of a message-passing network}
\outcome{21--26}{say what a topological order is, why there is usually more
  than one, and why they all give the same answer}
\outcome{27--34}{read a stationary distribution as an eigenvector, and say what
  a node with no outgoing edge does to one}
\end{outcomes}

\begin{quiz}
\item A graph has $\val{p13.g.ne}$ edges. What do its degrees add up to?
  \teachesat{4--5}
  \answerto{$\val{p13.g.degsum}$, twice the edges, because every edge is
  counted once at each of its two ends. It is Program~\ref{prog:P12}'s double
  counting in a new coat.}
\item $\val{p13.big.n}$ nodes and $\val{p13.big.m}$ edges. Roughly how many
  cells does the adjacency matrix have? \teachesat{9--10}
  \answerto{$\val{p13.big.cells}$ \dash{} a million squared \dash{} against
  $\val{p13.big.slots}$ slots for the list, a factor of
  $\val{p13.big.ratio}$. Real graphs are sparse and the matrix is stored only
  when it is going to be multiplied.}
\item What does the $(i,j)$ entry of $A^{2}$ count? \teachesat{14--15}
  \answerto{Walks of length two from $i$ to $j$. The sum
  $\sum_k A_{ik}A_{kj}$ has a term of $1$ for each intermediate node joined to
  both, which is Program~\ref{prog:P06}'s composition doing arithmetic.}
\item A directed graph with no cycle has how many topological orders?
  \teachesat{22--23}
  \answerto{Usually more than one \dash{} $\val{p13.topo.count}$ for the
  six-node graph in this program, out of $\val{p13.topo.total}$ orderings.
  What matters is that every one of them computes the same values.}
\item A page links nowhere. What does a random walk do when it lands on it?
  \teachesat{31--32}
  \answerto{It stops, and the probability mass leaves the graph: after
  $\val{p13.drain.zeroat}$ steps there is exactly none left. That is what the
  damping factor is for, and it is not a fudge.}
\end{quiz}

% ============================================================
\section{A set and a relation}
\label{sec:P13-what}
\index{graph!definition}

\begin{fr}
Program~\ref{prog:F10} said, in as many words, that a graph is a set plus a
relation and that saying so properly belonged here. This is that program, and
unlike the last one it does lean on Part~III: by section 3 the whole subject
will have turned into matrix multiplication.

Start concretely. Six services call one another: the gateway calls
\code{auth} and \code{search}; \code{search} calls \code{rank} and
\code{store}; \code{rank} calls \code{store}; and \code{auth}, \code{rank} and
\code{store} all write to \code{log}.

Forget the picture you are drawing. What is the smallest thing you could write
down that describes that completely?
\dotline
\nextframe
\end{fr}

\begin{fr}
\begin{ansblock}
Two sets: the services, and the pairs of them that are joined.
\end{ansblock}

That is the whole definition. A \textbf{graph} is a set $V$ of
\emph{vertices} and a set $E$ of pairs of them, its \emph{edges}. Here
$\lvert V \rvert = \val{p13.g.nv}$ and $\lvert E \rvert = \val{p13.g.ne}$.

Nothing about position, nothing about crossing lines, nothing about how it
looks. A drawing is one rendering of the relation and there are infinitely
many; the relation is the object.

\begin{note}
This is Program~\ref{prog:F10}'s point about sets arriving one level up. A set
answers \emph{is this one in it}; a graph answers \emph{are these two
joined}. Everything in this program is one of those two questions asked of a
bigger collection.
\end{note}

One thing the sentence above did leave out. \emph{The gateway calls
\code{auth}} is not the same statement as \emph{\code{auth} calls the
gateway}. What has to change in the definition to say so?
\nextframe
\end{fr}

\begin{fr}
\begin{ansblock}
The pairs have to be ordered.
\end{ansblock}

An \textbf{ordered} pair $(u, v)$ makes the graph \emph{directed}: the edge
goes from $u$ to $v$ and not back. An unordered pair $\{u, v\}$ makes it
undirected, and then the relation is symmetric by construction.

That single choice is the difference between a call graph, an import graph and
a dependency graph on one hand \dash{} all directed \dash{} and a friendship
graph or a similarity graph on the other. It is worth deciding explicitly,
because almost every graph an engineer meets is directed and almost every
picture drawn on a whiteboard forgets the arrows.

\mermaidfig{p13-set-plus-relation}{A graph is two sets, and the choice of
ordered or unordered pairs is the whole of what \enquote{directed} means.}{two
sets and a relation}
\end{fr}

\begin{fr}
Read the graph as undirected for a moment. A vertex's \textbf{degree} is the
number of edges at it, which here runs up to $\val{p13.g.degmax}$.

Add all six degrees together. Do it before reading on, and say what the total
has to be in terms of $\lvert E \rvert$.
\dotline
\nextframe
\end{fr}

\begin{fr}
\ans{$\val{p13.g.degsum}$, which is $2\lvert E \rvert$}
And the reason is one sentence: every edge has two ends, so it contributes $1$
to the degree at each of them, and adding the degrees counts every edge
exactly twice.

\textbf{That is Program~\ref{prog:P12}'s double counting in a new coat.} It is
usually called the handshake lemma, and the script checks it on
$\val{p13.hand.trials}$ random graphs \dash{} which, as in
Program~\ref{prog:P04}, is testing the code rather than looking for a
counterexample, because a counterexample would refute a proof.

\begin{note}
The immediate consequence is worth having: since the degrees sum to an even
number, \textbf{the number of vertices of odd degree is even.} You cannot have
a network in which exactly one machine has an odd number of links, and that
kind of parity argument is often the cheapest sanity check on a graph you have
just built.
\end{note}
\end{fr}

\begin{fr}
Now the other end of the range. A graph in which \emph{every} pair is joined
is called complete. If it has $\val{p13.complete.n}$ vertices, how many edges
does it have?
\nextframe
\end{fr}

\begin{fr}
\ans{$\val{p13.complete.e}$}
Which is $\binom{n}{2}$, and is Program~\ref{prog:F10}'s pair count and
Program~\ref{prog:P12}'s to the digit \dash{} because an undirected edge
\emph{is} an unordered pair, so counting edges in a complete graph is counting
pairs and nothing else.

\begin{note}
The script reads that program's committed pair count out of its values file
and fails the build if $\frac{n(n-1)}{2}$ does not reproduce it. Three
programs now quote one computation, and none of them can drift.
\end{note}

So the edges of any graph on $n$ vertices are a \emph{subset} of the
$\binom{n}{2}$ pairs, and the fraction actually present is its
\textbf{density}. That word is about to decide how the thing is stored.
\end{fr}

\begin{fr}
One place this is not an analogy. Program~\ref{prog:P12} deduplicated a corpus
by comparing every pair of documents. Make that a graph: the documents are the
vertices, and two are joined when their similarity clears a threshold.

Everything that program said about pairs is now a statement about edges. The
number of comparisons is the number of pairs, which is the complete graph; the
number of \emph{matches} is the number of edges, which is far smaller; and a
group of documents that are all duplicates of one another is a set of vertices
joined to each other and to nothing else.

\begin{aibox}
That last object has a name and you have met it without it. A
\textbf{connected component} is a set of vertices reachable from one another
and from nothing outside, and deduplicating a corpus is exactly the job of
finding the components and keeping one document from each.

Naming it matters because it changes what you look for. \enquote{Which pairs
match} is a question about $\binom{n}{2}$ things. \enquote{Which components
are there} is a question about $n$ of them, and there is a standard algorithm
for it that a graph library already has.
\end{aibox}
\end{fr}

% ============================================================
\section{Two encodings, two costs}
\label{sec:P13-encodings}
\index{graph!adjacency list}

\begin{fr}
There are two ways to write the relation down, and they are not
interchangeable.

The \textbf{adjacency matrix} $A$ is $n \times n$ with $A_{ij} = 1$ when there
is an edge from $i$ to $j$ and $0$ otherwise. The \textbf{adjacency list}
keeps, for each vertex, the list of vertices it points at.

Take a graph with $\val{p13.big.n}$ vertices and $\val{p13.big.m}$ edges,
which is an ordinary size for a citation network or a social graph. How many
cells does the matrix have?
\dotline
\nextframe
\end{fr}

\begin{fr}
\ans{$\val{p13.big.cells}$}
A million squared. At one byte a cell \dash{} which is already absurd, since a
cell holds one bit of information \dash{} that is
$\val{p13.big.tb}$ terabyte to store $\val{p13.big.m}$ edges.

The list needs one head per vertex and two ends per edge:
$\val{p13.big.slots}$ slots, a factor of $\val{p13.big.ratio}$ smaller.

\begin{trapbox}
\textbf{The matrix is not the definition and the list is not the
implementation.} Both are encodings of the same relation and neither is more
real than the other.

What makes the matrix look canonical is that it is what the mathematics is
written in, and what makes the list look practical is that it is what fits.
The honest statement is that they answer different questions cheaply, and the
density decides which question you are going to be asking. Here the density is
$\val{p13.big.density}$ \dash{} two edges in every hundred thousand possible
ones \dash{} and \textbf{almost every graph an engineer meets is sparse like
that}, which is why almost every graph library defaults to the list.
\end{trapbox}
\end{fr}

\begin{fr}
The operations differ as sharply as the memory does, and it is worth working
out which way round.

Two questions. \emph{Is there an edge from $u$ to $v$?} and \emph{what are all
of $u$'s neighbours?} For each of them, say which encoding answers faster and
why.
\dotline
\nextframe
\end{fr}

\begin{fr}
\begin{ansblock}
The matrix answers the first in one step and the second in $n$; the list
answers the second in as many steps as $u$ has neighbours, and the first in
the same.
\end{ansblock}

The matrix has the answer to \emph{is there an edge} at a known address, so it
is one lookup whatever the graph looks like. Asking it for the neighbours
means reading a whole row, most of which is zeros.

The list has the neighbours already gathered, so that question costs what the
answer is worth. Asking it \emph{is there an edge} means searching the list,
which costs the same as reading it.

\begin{note}
Program~\ref{prog:P03} is the program for saying this in $O$-notation, and
this one deliberately does not. \emph{This lookup reads a row of a million
cells} is a more useful sentence in a design review than \emph{this lookup is
linear}, and it is available without any notation at all.
\end{note}
\end{fr}

\begin{fr}
So the list wins on memory, wins on neighbours, and ties on lookup. The matrix
would be an odd thing to keep at all \dash{} except for one property that has
nothing to do with looking anything up.

\textbf{You can multiply a matrix.} The next section is what that buys, and it
is the reason the adjacency matrix survives in a subject where it is
$\val{p13.big.ratio}$ times too large to store.
\end{fr}

% ============================================================
\section{Walks are matrix powers}
\label{sec:P13-walks}
\index{graph!walk}

\begin{fr}
A \textbf{walk} of length $k$ is a sequence of $k$ edges laid end to end. It
may revisit a vertex; a walk that does not is a \emph{path}.

Now the question the whole section rests on. Write out the $(i,j)$ entry of
$A^{2}$ as a sum, using nothing but the definition of matrix multiplication:
\[ (A^{2})_{ij} = \sum_{k} A_{ik}A_{kj} \]
and say what each term of that sum is, given that every $A$ entry is $0$ or
$1$.
\dotline
\nextframe
\end{fr}

\begin{fr}
\begin{ansblock}
Each term is $1$ exactly when $i$ points at $k$ \emph{and} $k$ points at $j$,
and $0$ otherwise. So the sum counts the intermediate vertices \dash{} which
is to say, the walks of length two from $i$ to $j$.
\end{ansblock}

\textbf{$(A^{k})_{ij}$ is the number of walks of length $k$ from $i$ to $j$},
and the induction from two to $k$ is the same sentence again.

Notice where that came from. Program~\ref{prog:P06} argued that a matrix
product \emph{is} a composition and that row-meets-column is the arithmetic
which falls out of it. A walk of two steps is the composition of two steps,
so the walk count is a matrix product; the theorem is that argument read in
the other direction, and no new machinery was needed for it.

\begin{note}
The script checks it rather than deriving it twice: for every pair of vertices
and every length up to $\val{p13.walk.maxk}$ it enumerates the walks by brute
force and requires the count to equal the matrix entry \dash{}
$\val{p13.walk.checks}$ comparisons, all exact over the integers. An
enumeration cannot be wrong about what a walk is.
\end{note}
\end{fr}

\begin{fr}
Try it on the services. How many walks of length $\val{p13.walk.k}$ run from
the gateway to \code{log}?
\nextframe
\end{fr}

\begin{fr}
\ans{$\val{p13.walk.count}$}
The two routes through \code{search}: by way of \code{rank} and by way of
\code{store}. Reading them off $A^{3}$ takes no search and no traversal, and
the same multiplication answers the question for every pair at once.

Here is a second reading of the same theorem, worth a moment because it is
where people slip. On an \emph{undirected} graph, what is the diagonal entry
$(A^{2})_{ii}$?
\dotline
\nextframe
\end{fr}

\begin{fr}
\begin{ansblock}
The degree of $i$: the walks of length two from $i$ to itself are exactly
\enquote{step to a neighbour, step back}.
\end{ansblock}

\begin{trapbox}
\textbf{A walk is not a path, and the diagonal is where the difference shows.}
Going out and back is a perfectly good walk of length two and there are as
many of them as the vertex has neighbours; it is not a path, and it is
certainly not a cycle.

That is the trap in \enquote{$A^{k}$ counts the routes}, said too quickly. If
what you want is paths \dash{} routes that do not repeat a vertex \dash{}
matrix powers do not give them, and the counting problem is genuinely harder.
The powers answer the question they answer, which is the one about walks, and
that question is the one message passing turns out to ask.
\end{trapbox}

\mermaidfig{p13-walks-are-powers}{One step is the matrix, $k$ steps is its
$k$-th power, and message passing is that multiplication.}{steps, powers,
passing}
\end{fr}

\begin{fr}
That last sentence is the payoff, so here it is in full.

A graph neural network gives each vertex a feature vector, stacks them as the
rows of a matrix $X$, and one layer computes something of the shape
\[ X' = \sigma(AXW) \]
Read the three factors from the right: $W$ mixes each vertex's own features,
$A$ replaces each vertex's row by the sum of its neighbours' rows, and
$\sigma$ is the nonlinearity Program~\ref{prog:P06} showed the composition
would collapse without.

\textbf{Message passing is a multiplication by the adjacency matrix.} It is
not analogous to one and it is not implemented as one for convenience; the
sentence \enquote{each vertex aggregates from its neighbours} and the
expression $AX$ are the same statement.

So: after $k$ such layers, which vertices can have influenced a given
vertex's row?
\dotline
\nextframe
\end{fr}

\begin{fr}
\ans{Exactly those within $k$ steps of it}
Because $k$ layers apply $A$ $k$ times, and $A^{k}$'s pattern of non-zeros is
precisely the pairs joined by a walk of length $k$ \dash{} which is the
theorem two frames ago, arriving as a statement about architecture.

Measured on the services: at one hop the gateway reaches
$\val{p13.reach.one}$ vertices, and at $\val{p13.hops}$ hops it reaches all
$\val{p13.reach}$. The script checks that reach against a breadth-first
search, which knows nothing about matrices, for every vertex and every depth
up to $\val{p13.walk.maxk}$.

\begin{aibox}
\textbf{Depth in a message-passing network is reach, not capacity}, and that
is a real difference from every other architecture in this book. Adding a
layer to a feed-forward network gives it more to compute with; adding a layer
to a graph network gives every vertex a wider neighbourhood, whether or not
that was wanted.

So the question \enquote{how many layers} has an answer that comes from the
graph rather than from the budget: it is how far information genuinely has to
travel. On a graph of diameter $\val{p13.hops}$ a fourth layer adds nothing
new to reach, and on a small-world graph two or three layers already touch
most of the vertices.
\end{aibox}

\begin{warning}
The failure mode that follows \dash{} that with enough layers every vertex
aggregates from nearly the whole graph and the representations converge to one
another \dash{} is usually called over-smoothing, and \textbf{this book has not
measured it}. What is established here is the reach, which is a fact about
$A^{k}$ and needs no experiment. Whether and when the representations collapse
is a claim about trained networks, and it would need one.
\end{warning}
\end{fr}

% ============================================================
\section{A DAG, and the order that makes evaluation well defined}
\label{sec:P13-dag}
\index{graph!directed acyclic}

\begin{fr}
Back to the call graph, and to the property it has that has not been used yet:
you cannot get back to where you started.

A \textbf{cycle} is a walk that returns to its starting vertex without
retracing an edge. A directed graph with none is \textbf{acyclic}, and a
directed acyclic graph is a \emph{DAG}.

The clause about retracing is what makes the word usable, and it is worth a
moment. Every undirected graph with even one edge has a walk that leaves a
vertex and comes back \dash{} the \enquote{out and back} of two frames ago
\dash{} so without that clause nothing would be acyclic and the word would
mean nothing. On a directed graph the clause is doing less work, because
following an edge back is only possible if the edge back exists; \emph{that}
is the reason acyclicity is a statement people make about directed graphs and
almost never about undirected ones.

Suppose you have to run those six services in some order, each after
everything it depends on. What does acyclicity buy you?
\dotline
\nextframe
\end{fr}

\begin{fr}
\begin{ansblock}
That such an order exists at all.
\end{ansblock}

An order in which every vertex comes after everything pointing at it is a
\textbf{topological order}, and the theorem is that a directed graph has one
\emph{if and only if} it is acyclic. One direction is easy and worth doing in
your head: if there were a cycle, each of its vertices would have to come
after the next one round, and a finite list cannot do that.

The other direction is what makes the word \emph{acyclic} worth saying, and
this book does not prove it \dash{} Program~\ref{prog:P14} is where reading
and stating a theorem properly is taught.

Now count them. Our graph has $\val{p13.g.nv}$ vertices, so there are
$\val{p13.topo.total}$ orderings in all. How many are topological orders?
\nextframe
\end{fr}

\begin{fr}
\ans{$\val{p13.topo.count}$}
Four, out of $\val{p13.topo.total}$ \dash{} a fraction of
$\val{p13.topo.frac}$. The edges are a strong constraint and most orderings
break at least one of them.

The script finds them by testing every permutation, which is
Program~\ref{prog:P12}'s $n!$ arriving as a cost: it is fine at six and
hopeless at twenty, and a real scheduler uses an algorithm rather than a
search. What matters here is not how they are found but that there is more
than one.

Which raises the question this section exists for. Those
$\val{p13.topo.count}$ orders run the services in genuinely different
sequences. Do they compute the same thing?
\dotline
\nextframe
\end{fr}

\begin{fr}
\ans{Yes, always}
Give each service a value \dash{} say one more than the largest value among
the services that call it \dash{} and every topological order produces
identical values for all six. Not usually, not nearly: identically, and the
script asserts it over all $\val{p13.topo.count}$ of them.

\transcript{p13-order-and-answer}

The reason is short enough to hold. A vertex's value depends only on its
parents; a topological order computes every parent before the vertex; so
whichever order is chosen, each vertex sees the same parent values when its
turn comes. Nothing else about the order can reach it.

\mermaidfig{p13-order-not-answer}{Many orders, one answer \dash{} which is what
lets a scheduler pick whichever it likes.}{order is free, answer is not}
\end{fr}

\begin{fr}
\begin{aibox}
That is the whole reason four unrelated-looking things are one object.

A \textbf{build system} has files depending on files, runs them in a
topological order, and is free to parallelise the parts that do not depend on
one another \dash{} which is exactly the freedom the
$\val{p13.topo.count}$ orders represent. An \textbf{agent workflow} has steps
depending on steps. A \textbf{neural network's forward pass} has tensors
depending on tensors. A \textbf{dataflow query plan} has operators depending
on operators.

All four are DAGs, all four are evaluated in a topological order, and all four
are deterministic for the same reason. When a framework calls the thing it
builds a \emph{computation graph}, it is not a metaphor.
\end{aibox}

Program~\ref{prog:P16} is where that graph is used rather than defined: the
backward pass walks the same DAG in the reverse of a topological order, which
is what makes \code{loss.backward()} well defined and is why the activations
have to still be there when it arrives.
\end{fr}

\begin{fr}
And the negative case earns its own frame, because it is the one people meet.

Add one edge to the call graph \dash{} let \code{log} call the gateway \dash{}
and ask for the topological orders again. The transcript above already shows
the answer: the empty list.

\begin{trapbox}
\textbf{A cycle does not make evaluation hard, it makes it undefined.} There is
no order to run the services in, because each would have to run after the
next. A build system reports a dependency cycle and stops; it is not being
unhelpful, there is genuinely nothing to do.

The two ways out are both worth recognising. \emph{Break the cycle} \dash{}
find the edge that should not be there, which in a call graph is usually a
layering violation. Or \emph{unroll it}: a recurrent network is a cycle in the
architecture and a DAG in the computation, because the same weights are
applied at each of a fixed number of steps and those steps are distinct
vertices. That is why a recurrent network can be differentiated at all, and
why its memory grows with the number of steps unrolled.
\end{trapbox}
\end{fr}

% ============================================================
\section{Where a random walk settles}
\label{sec:P13-walks-random}
\index{graph!random walk}
\index{graph!PageRank}

\begin{fr}
One last object, and it puts Part~III's machinery on a graph.

Instead of asking which vertices are reachable, follow the edges at random:
from a vertex, pick one of its outgoing edges uniformly and go. The matrix of
those probabilities is $P$, where $P_{ij}$ is the chance of stepping from $i$
to $j$ \dash{} the adjacency matrix with each row divided by its own row sum.

What must each row of $P$ add up to, and why?
\dotline
\nextframe
\end{fr}

\begin{fr}
\ans{$1$, because the walk goes somewhere}
Each row is a probability distribution over where the next step lands, and
dividing a row of ones by the number of ones is exactly what makes it one.

Now put a distribution $p$ over the vertices \dash{} the chance of being at
each \dash{} and take a step. The new distribution is $P\T p$: the chance of
arriving at $j$ is the sum over $i$ of \emph{being at $i$} times
\emph{stepping from $i$ to $j$}, which is a matrix-vector product and nothing
more.

Start uniform and iterate. Say what you expect to happen to $p$.
\nextframe
\end{fr}

\begin{fr}
\begin{ansblock}
It settles: after enough steps the distribution stops changing.
\end{ansblock}

Measured on a four-page link graph, after $\val{p13.pi.steps}$ steps:
$\val{p13.pi.0}$, $\val{p13.pi.1}$, $\val{p13.pi.2}$, $\val{p13.pi.3}$. The
exact values are
$\frac{\val{p13.exact.num.0}}{\val{p13.exact.den.0}}$,
$\frac{\val{p13.exact.num.1}}{\val{p13.exact.den.1}}$,
$\frac{\val{p13.exact.num.2}}{\val{p13.exact.den.2}}$ and
$\frac{\val{p13.exact.num.3}}{\val{p13.exact.den.3}}$, solved over the
rationals rather than iterated, and the iteration agrees with them to better
than $10^{-\val{p13.pi.bound}}$.

A distribution that a step leaves alone is called \textbf{stationary}. Write
that condition down as an equation in $P\T$ and $p$, and say what you are
looking at.
\dotline
\nextframe
\end{fr}

\begin{fr}
\ans{$P\T p = p$, which is an eigenvector for eigenvalue $1$}
Program~\ref{prog:P10} defined exactly this: a direction the matrix does not
turn, only stretches, and here the stretch factor is $1$.

So the stationary distribution of a random walk is an eigenvector calculation,
and everything that program said applies unchanged \dash{} including that
finding it by repeated multiplication is what power iteration is, and that the
rate it settles at is governed by the second eigenvalue.

\begin{note}
The script does not iterate and hope. It solves $(P\T - I)p = 0$ with
$\sum p = 1$ exactly over the rationals, by the same Gaussian elimination
Program~\ref{prog:P04} wrote and Programs~\ref{prog:P08}
and~\ref{prog:P09} reused, and then asserts $P\T p = p$ with no tolerance at
all. \enquote{Stationary} is a claim about equality, and a float comparison
would have made it a claim about a threshold.
\end{note}
\end{fr}

\begin{fr}
That is \textbf{PageRank}, and the definition is the whole of it: a page's
rank is its share of the time an endless random walk spends there. A page is
important because important pages link to it, which sounds circular and is
instead an eigenvector equation.

That distinction is worth having in words, because \enquote{circular} is the
first objection anybody raises. A circular \emph{definition} explains a term
using itself and explains nothing. What $P\T p = p$ states is not a definition
but a \emph{condition}: it says nothing about what the ranking is, and instead
demands that whatever it is, it agrees with itself after one step of the walk.
Whether any such ranking exists at all is then a genuine question with a
genuine answer, and the answer is the eigenvector \dash{} which exists, and
under the condition of the next two frames is unique.

But the walk as described breaks on a graph the web certainly contains. What
happens when it lands on a page with no outgoing links at all?
\dotline
\nextframe
\end{fr}

\begin{fr}
\begin{ansblock}
It has nowhere to go, and the probability mass leaves the graph.
\end{ansblock}

That row of $P$ is all zeros, so it is not a distribution and the iteration
does not preserve the total. Measured on a three-page graph with one dead end:
the mass left after each step is
$\frac{\val{p13.drain.num1}}{\val{p13.drain.den1}}$, then
$\frac{\val{p13.drain.num2}}{\val{p13.drain.den2}}$, then
\textbf{exactly zero} by step $\val{p13.drain.zeroat}$.

Not nearly zero. Zero, over the rationals, with everything gone.
\end{fr}

\begin{fr}
The repair is one line and it is usually described as a hack, which it is not.

With probability $\val{p13.damp}$ follow a link as before; with the remaining
probability jump to a vertex chosen uniformly at random. A page with no links
jumps every time. Two things follow at once, and both are the point.

\textbf{Every entry of the matrix is now positive}, so every page is reachable
from every other in one step, and the walk cannot get trapped in a subgraph
that links only to itself. And \textbf{nothing drains}, because every row sums
to one again. Measured on the same three-page graph, the damped walk settles
at $\val{p13.damped.0}$, $\val{p13.damped.1}$ and $\val{p13.damped.2}$, and
the dead end keeps the largest share rather than swallowing everything.

\begin{note}
The value $\val{p13.damp}$ is conventional rather than derived, and this book
will not pretend otherwise. What \emph{is} derived is that the damping has to
be there: without it the fixed point either fails to exist or fails to be
unique, and with it the matrix is positive, which is the condition the
relevant theorem asks for.
\end{note}
\end{fr}

\begin{fr}
One closing observation, which is what this program has really been about.

Three questions were asked of one object and each turned into arithmetic
already in the book. \emph{What can reach what} became a power of a matrix.
\emph{In what order may this be evaluated} became a property that costs
nothing to check and buys determinism. \emph{Where does attention settle}
became an eigenvector.

\textbf{None of the three needed a picture}, and the picture is the thing most
people reach for first. A graph drawn on a whiteboard is a rendering of a
relation; the relation is what the arithmetic is about, and it is the only
part that survives being a million vertices wide.
\end{fr}

% ============================================================
\begin{summarybox}
\sumitem{1--3}{\result{\textbf{A graph is a set of vertices and a set of pairs of
  them}, and nothing else \dash{} no position, no drawing.
  \textbf{Ordered pairs make it directed}, which is what a call graph, an
  import graph and a dependency graph have and a whiteboard sketch usually
  loses}.}
\sumitem{4--5}{\result{\textbf{The degrees sum to twice the edges}, because every edge
  is counted at each of its two ends: $\val{p13.g.degsum}$ for
  $\val{p13.g.ne}$ edges. So the number of odd-degree vertices is even}.}
\sumitem{6--8}{\result{A complete graph on $\val{p13.complete.n}$ vertices has
  $\val{p13.complete.e}$ edges, which is Programs~\ref{prog:F10}
  and~\ref{prog:P12}'s pair count \dash{} an undirected edge \emph{is} an
  unordered pair. Deduplication is finding connected components}.}
\sumitem{9--10}{\result{\textbf{Matrix or list is a question about density and about
  which operation you repeat.} At $\val{p13.big.n}$ vertices and
  $\val{p13.big.m}$ edges the matrix is $\val{p13.big.cells}$ cells against
  $\val{p13.big.slots}$ slots, and the density is
  $\val{p13.big.density}$}.}
\sumitem{11--13}{\result{The matrix answers \emph{is there an edge} in one step and
  the list answers \emph{who are the neighbours} in as many steps as there are
  neighbours. \textbf{The matrix earns its place by being multipliable}, not
  by being lookupable}.}
\sumitem{14--15}{\result{\textbf{$(A^{k})_{ij}$ counts the walks of length $k$ from $i$
  to $j$}, because each term of $\sum_k A_{ik}A_{kj}$ is $1$ exactly when the
  two steps both exist \dash{} Program~\ref{prog:P06}'s composition doing
  arithmetic}.}
\sumitem{16--18}{\result{\textbf{A walk is not a path}: $(A^{2})_{ii}$ is the degree of
  $i$, counting \enquote{out and back}. Counting paths is a harder problem and
  matrix powers do not do it}.}
\sumitem{19--20}{\result{\textbf{Message passing is a multiplication by the adjacency
  matrix}, so $k$ layers reach exactly $k$ hops: $\val{p13.reach.one}$
  vertices at one hop here and all $\val{p13.reach}$ at $\val{p13.hops}$.
  \textbf{Depth is reach, not capacity.}}}
\sumitem{21--23}{\result{\textbf{A directed graph has a topological order if and only
  if it is acyclic}, and there is usually more than one:
  $\val{p13.topo.count}$ here out of $\val{p13.topo.total}$}.}
\sumitem{24--25}{\result{\textbf{Every topological order computes the same values},
  because each vertex depends only on parents that all of them compute first.
  That is why a build system, an agent workflow, a query plan and a forward
  pass are one object, and why a scheduler may parallelise freely}.}
\sumitem{26}{\result{\textbf{A cycle makes evaluation undefined rather than hard.}
  Break it, or unroll it \dash{} which is what a recurrent network is, and why
  its memory grows with the steps unrolled}.}
\sumitem{27--30}{\result{A random walk's step is $P\T p$, and
  \textbf{a stationary distribution is an eigenvector for eigenvalue $1$}
  \dash{} Program~\ref{prog:P10} under another name. That is PageRank, and it
  is checked here as an exact equality rather than a tolerance}.}
\sumitem{31--34}{\result{\textbf{A vertex with no outgoing edge drains the walk}:
  exactly no mass is left after $\val{p13.drain.zeroat}$ steps. Damping at
  $\val{p13.damp}$ makes every entry positive and every row sum to one, which
  is what buys a unique fixed point \dash{} the constant is conventional, the
  need for it is not}.}
\end{summarybox}

\canyou

\begin{testexercises}
\item A directed graph has $12$ vertices and $30$ edges. What do the
  out-degrees sum to? And if it were undirected, the degrees?
  \answerto{$30$ either way for the out-degrees \dash{} each directed edge
  leaves exactly one vertex. Undirected it is $60$, twice the edges, because
  each edge is counted at both ends. The two answers differ by a factor of two
  and confusing them is the commonest slip in this subject.}
\item You have $50\,000$ vertices and $200\,000$ edges. Give the matrix and
  list sizes, and say which you would store.
  \answerto{$50\,000^{2} = \num{2.5e9}$ cells against
  $50\,000 + 400\,000 = \num{450000}$ slots, a factor of about $5500$. Store
  the list; the density is $\num{1.6e-4}$. Keep a sparse matrix as well only
  if you intend to multiply by it.}
\item $A$ is the adjacency matrix of an undirected graph. What is the sum of
  all the entries of $A$? And of $A^{2}$'s diagonal?
  \answerto{Both are twice the number of edges. The first because each edge
  sets two entries; the second because $(A^{2})_{ii}$ is the degree of $i$ and
  the degrees sum to $2\lvert E \rvert$. The same quantity by two routes,
  which is the handshake lemma again.}
\item A build has $8$ targets and one of them, transitively, depends on
  itself. What will the build system do, and what are your two options?
  \answerto{It reports a dependency cycle and refuses to run, because there is
  no order in which each target follows what it needs. Either break the cycle
  \dash{} usually a layering violation \dash{} or unroll it into distinct
  copies, which is what a recurrent network does and why its memory grows with
  the number of steps.}
\item A graph neural network has $6$ layers and its graph has diameter $3$.
  What does the sixth layer add to any vertex's reach?
  \answerto{Nothing. After $3$ layers every vertex already aggregates from the
  whole connected component, so layers four to six add parameters and
  computation and no new information. Depth in a message-passing network is
  reach, and reach stops at the diameter.}
\end{testexercises}

\begin{furtherproblems}
\item Show that the number of vertices of odd degree is even.
  \answerto{The degrees sum to $2\lvert E \rvert$, which is even. The
  even-degree vertices contribute an even total, so the odd-degree ones must
  too \dash{} and a sum of odd numbers is even only when there is an even
  number of them.}
\item $A$ has a row of zeros. What does that say about the graph, and what
  does it do to $A^{k}$?
  \answerto{That vertex has no outgoing edge. Row $i$ of $A^{k}$ is then zero
  for every $k$, so no walk of any length leaves it \dash{} which is the same
  dead end that drains a random walk in section 5, seen in the adjacency
  matrix rather than in the transition matrix.}
\item Why does counting paths not reduce to a matrix power the way counting
  walks does?
  \answerto{Because \enquote{does not repeat a vertex} is a condition on the
  whole sequence, and matrix multiplication combines steps without knowing
  what came before them. The sum $\sum_k A_{ik}A_{kj}$ cannot exclude the
  terms where $k$ equals $i$ or $j$ without being told, and at length three
  the condition is about the whole route rather than one term.}
\item A DAG has $n$ vertices and no edges at all. How many topological orders?
  And with a single chain of $n-1$ edges?
  \answerto{$n!$ with no edges \dash{} every ordering qualifies \dash{} and
  exactly $1$ for a chain, which fixes the order completely. Every DAG sits
  between those two, and the count is a measure of how much freedom a
  scheduler has.}
\item The damping factor is often described as \enquote{the chance a surfer
  gets bored}. What does it actually buy, mathematically?
  \answerto{Two things. Every entry of the matrix becomes positive, so no set
  of pages can trap the walk, which is what makes the fixed point unique; and
  every row sums to one even for a page with no links, so nothing drains away.
  The story about boredom motivates it; those two properties are what the
  theorem needs.}
\end{furtherproblems}
